\subsection{Conceptualisations for Design Approach}
\subsubsection{Contactless Capacitive Water Level Sensing}
The strip will be placed vertically on the exterior side of the brita/pitcher to detect exact water level. 
Works as long as plastic is not too thick which britas fit the description. \\

\noindent\textbf{Pros}
    \begin{itemize}
    \item Low-cost
    \item No contact
    \item Small mechanical approach
    \item Customizable for future addons
    \end{itemize}
\noindent\textbf{Cons}
    \begin{itemize}
    \item Need a good environment for sensor, not humid so not great inside a fridge
    \end{itemize}

\subsubsection{Weight based, Load}
The design places the pitcher unto a load cell to measure the weight. We will calibrate it to the changes in weight.\\
\noindent\textbf{Pros}
    \begin{itemize}
    \item Highest accuracy of all designs
    \item Fits to many pitcher dimensions
    \end{itemize}
\noindent\textbf{Cons}
    \begin{itemize}
    \item Need a good environment for sensor, not humid so not great inside a fridge
    \end{itemize}

\subsubsection{Ultrasonic}
An ultrasonic sensor estimates the water level by measuring the echo time of sound waves reflected from the water surface.\\
\noindent\textbf{Pros}
    \begin{itemize}
    \item Inexpensive
    \item Easy implementation
    \end{itemize}
\noindent\textbf{Cons}
    \begin{itemize}
    \item Humidity problem
    \end{itemize}

\subsubsection{Laser}
A laser or time-of-flight sensor measures the distance from the sensor to the water surface.
The measured distance is converted into liquid height and volume. \\
\noindent\textbf{Pros}
    \begin{itemize}
    \item Continous measurements
    \item Non contact
    \end{itemize}
\noindent\textbf{Cons}
    \begin{itemize}
    \item Humidity problem
    \item Clear line of sight
    \end{itemize}
